\chapter{Atomhülle}

\section{2024-11-27}

\startframedtipp
Für die Abi-Klausur:

Planksche konstante anhand LED bestimmen sehr beliebte Abi Aufgabe.
\stopframedtipp

\startitemize[packed]
\item inverser Photoeffekte : Photonen (Licht) in Portionen mit {\em bestimmter} Energie
\item es stehen nicht alle mathematisch möglichen Energiewerte zur Verfügung
\item Tritt das auch wo anders auf?
\stopitemize

\subsection{Notizen zu uns bekannten Atommodellen}

\startitemize[packed]
\item Demokrit (antikes Griechenland, ca 400 v. Chr.):
\startitemize[packed]
\item "altomos" (Das Kleinste), alles besteht aus kleinsten Teilchen
\stopitemize
\item Dalton: Atome sind kleinste Teilchen, unterscheiden sich nach Element
\item Thomson: Atome aus Elektronen (negativ geladen, Rosinen) in festen Rümpfen (positiv, Kuchen); Rosinenkuchenodell
\item Rutherford: (Goldfolie mit $\alpha$-Strahlung) $e^-$ außen, Kern ($t$) innen, sonst leer
\item Bohr: $e^-$ kreisen auf festen Bahnen mit zugehöriger Energie, Kern in Mitte
\item Schalenmodell: $e^-$ mit fester Energie, Position und Bewegung nicht bekannt
\item Orbitalmodell: $e^-$ mit Energie und Aufenthaltswahrscheinlcihkeiten für Bereiche
% \item Bohrsches Atommodell
% \startitemize[packed]
% \item 1913
% \item schwere, positiv geladener Atomkern
% \item leichte, negativ geladenen Elektronen
% \stopitemize
% \item Thomsonsches Atommodell
% \startitemize[packed]
% \item 1903
% \item Das Atom besteht aus gleichmäßig verteielter, positiv geladener Masse
% \stopitemize
\stopitemize


\subsection{Herleitung}

\startitemize[packed]
\item Beginnt mit (korrekter) Annahme bzw. Gleichung, Gültigkeit wird per Text begründet
\item Äquivalenzumformung
\item Zusatzannahme per Text erklärung
\item Abschlussatz
\stopitemize

\subsection{Geschwindigkeit Elektronen im homogenen elektrischen Feld}

\startframedtask
Situation: $e^-$ wird freigesetzt und beschleunigt.

Wir gehen davon aus, dass die elektrische Energie im homogenen elektrischen Feld vollständig in kinetische Energie umgewandelt
\stopframedtask

\startplaceformula[formel:autom:energie:el:kin]
\startformula
\startalign[n=4]
\NC E_{el} \NC = E_{kin} \NC \NC \NR
\NC e \cdot U = q \cdot U \NC = \frac{1}{2} mv^2 \NC \quad \| \cdot 2 \NC \NR
\NC 2eU \NC = m \cdot v^2 \NC \quad \| \div m \NC \NR
\NC \frac{2eU}{m} \NC = v^2 \NC \quad \| \sqrt{} \NC \NR
\NC \sqrt{\frac{2eU}{m}}\NC =v  \NC \NC \NR
\stopalign
\stopformula
\stopplaceformula

Damit ist die die Gleichung nach $v$ hergeleitet.


